% ID: SEL-D0043-OB-E04
% ESTADO: cotejada. Segunda pasada visual y borrador compilado revisados.
% Fuente: Selectividad/sel_2014_matematicas_aplicadas/EXAMEN -ANDAL-2.pdf
% SHA-256: fc025a25faa338cbb9d8d288443f0322dfbcc0464f7474a16a84e8dfe47f81d1
% Página: 2; opción B; ejercicio original 4.

1) En un centro docente la tercera parte de los alumnos estudia el idioma A, la mitad el idioma B y el resto el idioma C (cada alumno estudia sólo uno de estos idiomas).
\begin{enumerate}[label=\alph*)]
\item \textbf{(0.75 puntos)} Se desea seleccionar una muestra de 60 alumnos, mediante muestreo aleatorio estratificado con afijación proporcional al número de los alumnos de cada idioma. \textquestiondown{}Cómo debería estar conformada la muestra?
\item \textbf{(0.75 puntos)} En otra muestra seleccionada por el procedimiento anterior, el número de alumnos tomados del idioma A es 14. Determine cuántos se han elegido de los otros dos idiomas.
\end{enumerate}
2) \textbf{(1 punto)} Una población tiene 5 elementos. Mediante muestreo aleatorio simple se seleccionan muestras de tamaño 3, siendo la desviación típica de sus medias 2 y la media de las medias muestrales 7. \textquestiondown{}Cuánto valen la media y la varianza de la población?
